\section{Statement and height setup}
\label{mumford-AppII-statement}
\dcref{M:appII:II}

\begin{convention}[Arithmetic setting]
\label{mumford-conv-mordell-weil}
\dcref{M:appII:II}
Let $K$ be a number field and let $A/K$ be an abelian variety.  Write
$A(K)$ for its group of $K$-rational points and fix a symmetric ample line
bundle $L$ on $A$.
\notready
\end{convention}

\begin{theorem}[Mordell--Weil]
\label{mumford-thm-mordell-weil}
The group $A(K)$ is finitely generated.
\dcref{M:appII:II}
\uses{mumford-conv-mordell-weil,mumford-thm-weak-mordell-weil,
 mumford-thm-mordell-weil-descent}
\notready
\end{theorem}

\begin{proposition}[Quadratic canonical height]
\label{mumford-prop-canonical-height}
There is a function $\hat h_L:A(\overline K)\to\R$ and a symmetric bilinear
pairing $\langle\ ,\ \rangle_L$ on $A(\overline K)$ modulo torsion such that
\[
 \hat h_L(nx)=n^2\hat h_L(x),\qquad
 \langle x,y\rangle_L=\tfrac12(\hat h_L(x+y)-\hat h_L(x)-\hat h_L(y)),
\]
and $\hat h_L(x)\ge0$, with equality precisely for torsion points.  It differs
from a Weil height attached to $L$ by a bounded function.
\dcref{M:appII:II}
\uses{mumford-conv-mordell-weil,mumford-lem-height-parallelogram}
\notready
\end{proposition}

\begin{lemma}[Parallelogram identity]
\label{mumford-lem-height-parallelogram}
\dcref{M:appII:II}
For a symmetric line bundle and a Weil height $h_L$, the cube identity yields
\[
 h_L(x+y)+h_L(x-y)-2h_L(x)-2h_L(y)=O(1)
\]
uniformly in $x,y$.  The sequence $4^{-r}h_L(2^rx)$ is then Cauchy; its limit
defines the canonical height, and polarization of the defect gives the
associated bilinear pairing.
\uses{mumford-thm-cube-schemes}
\notready
\end{lemma}

\section{Weak Mordell--Weil theorem}
\label{mumford-AppII-weak}
\dcref{M:appII:II}

\begin{proposition}[Kummer injection]
\label{mumford-prop-kummer-injection}
\dcref{M:appII:II}
For every integer $n>1$, the Kummer map gives an injection
\[
 A(K)/nA(K)\hookrightarrow H^1(K,A[n]).
\]
Its image is described by local divisibility conditions at the places of $K$.
\uses{mumford-prop-multiplication-n,mumford-def-cartier-dual}
\notready
\end{proposition}

\begin{lemma}[Finite Kummer quotient]
\label{mumford-lem-finite-kummer-quotient}
\dcref{M:appII:II}
For fixed $n$, the subgroup of $H^1(K,A[n])$ satisfying the unramified and
integrality conditions outside a finite set of places is finite.  Hence
$A(K)/nA(K)$ is finite.
\uses{mumford-prop-kummer-injection}
\notready
\end{lemma}

\begin{theorem}[Weak Mordell--Weil]
\label{mumford-thm-weak-mordell-weil}
For every integer $n>1$, $A(K)/nA(K)$ is finite.
\dcref{M:appII:II}
\uses{mumford-lem-finite-kummer-quotient}
\notready
\end{theorem}

\section{Descent and finite generation}
\label{mumford-AppII-descent}
\dcref{M:appII:II}

\begin{lemma}[Northcott property]
\label{mumford-lem-northcott-abelian}
\dcref{M:appII:II}
For a projective embedding supplied by a symmetric ample line bundle, points
of bounded degree over $K$ and bounded Weil height form a finite set.  The same
holds for the canonical height because it differs from a Weil height by a
bounded amount.
\uses{mumford-prop-canonical-height,mumford-thm-high-power-very-ample}
\notready
\end{lemma}

\begin{proposition}[Finite descent set]
\label{mumford-prop-finite-descent-set}
\dcref{M:appII:II}
Choose $n>1$ and representatives $x_1,\ldots,x_r$ for $A(K)/nA(K)$.  There is a
constant $C$ such that every class has a representative $x_i$ with
$\hat h_L(x_i)\le C$; subtracting suitable multiples of $n$ decreases the
height unless the point lies in a finite Northcott set.
\uses{mumford-thm-weak-mordell-weil,mumford-lem-northcott-abelian}
\notready
\end{proposition}

\begin{lemma}[Finite rational torsion]
\label{mumford-lem-rational-torsion-finite}
\dcref{M:appII:II}
The torsion points in $A(K)$ have canonical height zero and therefore lie in a
bounded-height set of bounded degree.  Northcott's theorem makes this set
finite, so $A(K)_{\mathrm{tors}}$ is finite.
\uses{mumford-prop-canonical-height,mumford-lem-northcott-abelian}
\notready
\end{lemma}

\begin{theorem}[Descent conclusion]
\label{mumford-thm-mordell-weil-descent}
\dcref{M:appII:II}
The finite descent set together with the finite torsion subgroup generates
$A(K)$.  Therefore $A(K)$ is a finitely generated abelian group.
\uses{mumford-prop-finite-descent-set,mumford-lem-rational-torsion-finite,
 mumford-prop-canonical-height}
\begin{proof}
Let $\Gamma$ be the subgroup generated by the finite set of representatives
and $A(K)_{\mathrm{tors}}$.  For any $x\in A(K)$, repeatedly replace $x$ by a
representative of its class modulo $nA(K)$ and divide the remaining difference
by $n$.  The quadratic height relation shows that the height of the remainder
strictly decreases unless it is in the bounded finite set from Northcott.
Induction on the height therefore expresses $x$ in $\Gamma$.
\end{proof}
\notready
\end{theorem}

\begin{corollary}[Finite torsion]
\label{mumford-cor-finite-torsion}
\dcref{M:appII:II}
The torsion subgroup $A(K)_{\mathrm{tors}}$ is finite.
\uses{mumford-lem-rational-torsion-finite}
\notready
\end{corollary}

\begin{remark}[Mordell--Weil proof interface]
\label{mumford-app2-mordell-weil-theorem}
The finite-generation theorem combines the weak quotient, canonical height, and
finite descent interfaces.
\dcref{M:appII:II}
\uses{mumford-thm-mordell-weil,mumford-app2-final-descent}
\notready
\end{remark}

\section{Supporting arithmetic interfaces}
\label{mumford-AppII-coverage-register}
The arithmetic interfaces in this register make the weak theorem, height
construction, and descent dependencies explicit.  Detailed proofs are treated
as external inputs where necessary.

\begin{proposition}[Finite weak quotient target]
\label{mumford-app2-weak-quotient-finite}
For each $n>1$, finiteness of $A(K)/nA(K)$ is the weak Mordell--Weil target used
by descent.
\dcref{M:appII:II}
\uses{mumford-thm-weak-mordell-weil}
\notready
\end{proposition}

\begin{proposition}[Height pairing target]
\label{mumford-app2-height-pairing-target}
The canonical height is quadratic, nonnegative for a symmetric ample bundle,
and positive modulo torsion on $A(K)$.
\dcref{M:appII:II}
\uses{mumford-prop-canonical-height,mumford-lem-rational-torsion-finite}
\notready
\end{proposition}

\begin{proposition}[Abstract descent criterion]
\label{mumford-app2-abstract-descent}
An abelian group with finite quotient by $n$ and a positive quadratic height with
finite bounded sets is finitely generated.
\dcref{M:appII:II}
\uses{mumford-thm-weak-mordell-weil,mumford-app2-height-pairing-target,
 mumford-lem-northcott-abelian}
\notready
\end{proposition}

\begin{lemma}[Division points and Kummer cocycle]
\label{mumford-app2-kummer-fundamental-lemma}
After adjoining an $n$-division point of $x\in A(K)$, the differences
$\sigma(z)-z$ define a cocycle with values in $A[n]$, independent of the
chosen division point up to coboundary.
\dcref{M:appII:II}
\uses{mumford-prop-kummer-injection}
\notready
\end{lemma}

\begin{lemma}[Finite division kernel]
\label{mumford-app2-finite-division-kernel}
The kernel of multiplication by $n$ is finite, and the Kummer cocycle vanishes
exactly when $x$ is divisible by $n$ in $A(K)$.
\dcref{M:appII:II}
\uses{mumford-prop-multiplication-n,mumford-app2-kummer-fundamental-lemma}
\notready
\end{lemma}

\begin{lemma}[Good model over an arithmetic open]
\label{mumford-app2-model-over-open}
After removing finitely many primes from the ring of integers of $K$, $A$ and
its ample bundle extend to a proper model with smooth group fibres.
\dcref{M:appII:II}
\uses{mumford-conv-mordell-weil}
\notready
\end{lemma}

\begin{corollary}[Etale division away from $n$]
\label{mumford-app2-etale-division-corollary}
On the good model, multiplication by $n$ is etale outside the finite set of
places dividing $n$ and the bad-reduction set.
\dcref{M:appII:II}
\uses{mumford-app2-model-over-open,mumford-prop-multiplication-n}
\notready
\end{corollary}

\begin{lemma}[Unramified Kummer classes]
\label{mumford-app2-arithmetic-unramified-classes}
Kummer classes unramified outside a fixed finite set of places form a finite
subgroup of $H^1(K,A[n])$; this uses finiteness of unit and ideal-class
quotients.
\dcref{M:appII:II}
\uses{mumford-app2-etale-division-corollary,mumford-prop-kummer-injection}
\notready
\end{lemma}

\begin{definition}[Normalized absolute values]
\label{mumford-app2-normalized-absolute-values}
Choose normalized absolute values on $K$ so that the product formula holds, and
define the projective height by the weighted sum of local logarithmic maxima.
\dcref{M:appII:II}
\notready
\end{definition}

\begin{proposition}[Projective height properties]
\label{mumford-app2-projective-height-properties}
The projective height is invariant under scaling, functorial under finite field
extension, and satisfies the standard bounded-error inequalities for products
and morphisms.
\dcref{M:appII:II}
\uses{mumford-app2-normalized-absolute-values}
\notready
\end{proposition}

\begin{proposition}[Northcott in projective space]
\label{mumford-app2-northcott-proposition}
Points of bounded degree and bounded projective height over a number field form
a finite set.
\dcref{M:appII:II}
\uses{mumford-app2-projective-height-properties}
\notready
\end{proposition}

\begin{definition}[Height of a line bundle]
\label{mumford-app2-line-bundle-height}
An embedding associated with an ample line bundle $L$ pulls the projective
height back to a Weil height $h_L$ on $A(\overline K)$.
\dcref{M:appII:II}
\uses{mumford-app2-projective-height-properties,mumford-thm-high-power-very-ample}
\notready
\end{definition}

\begin{proposition}[Comparison of equivalent heights]
\label{mumford-app2-height-comparison-proposition}
Heights attached to linearly equivalent bundles differ by a bounded function;
tensor products add heights up to bounded error.
\dcref{M:appII:II}
\uses{mumford-app2-line-bundle-height}
\notready
\end{proposition}

\begin{lemma}[Approximate quadraticity]
\label{mumford-app2-approximate-quadratic-lemma}
For symmetric $L$, the cube relation gives a uniform bounded defect in the
parallelogram identity for $h_L$.
\dcref{M:appII:II}
\uses{mumford-app2-height-comparison-proposition,mumford-thm-cube-schemes}
\notready
\end{lemma}

\begin{theorem}[Tate height construction]
\label{mumford-app2-tate-height}
The limit of the normalized heights $4^{-r}h_L(2^rx)$ exists and is a quadratic
height satisfying $\widehat h_L(nx)=n^2\widehat h_L(x)$.
\dcref{M:appII:II}
\uses{mumford-app2-approximate-quadratic-lemma,mumford-prop-canonical-height}
\notready
\end{theorem}

\begin{proposition}[Height positivity and finiteness]
\label{mumford-app2-height-positivity-finiteness}
For rational points, the canonical height is positive away from torsion, and
Northcott makes every bounded-height subset finite.
\dcref{M:appII:II}
\uses{mumford-app2-tate-height,mumford-app2-northcott-proposition}
\notready
\end{proposition}

\begin{remark}[Final descent step]
\label{mumford-app2-final-descent}
Choose representatives for the finite quotient modulo $n$ and repeatedly divide
the remainder.  The quadratic height decreases outside the finite Northcott
set, so induction gives a finite generating set.
\dcref{M:appII:II}
\uses{mumford-app2-abstract-descent,mumford-thm-mordell-weil-descent}
\notready
\end{remark}
